
\documentclass{article}
\usepackage{colortbl}
\usepackage{makecell}
\usepackage{multirow}
\usepackage{supertabular}

\begin{document}

\newcounter{utterance}

\centering \large Interaction Transcript for game `function\_detective', experiment `two\_numbers\_passive\_labeled\_pairs', episode 279 with Qwen3.6{-}35B{-}A3B{-}FP8{-}without{-}reasoning.
\vspace{24pt}

{ \footnotesize  \setcounter{utterance}{1}
\setlength{\tabcolsep}{0pt}
\begin{supertabular}{c@{$\;$}|p{.15\linewidth}@{}p{.15\linewidth}p{.15\linewidth}p{.15\linewidth}p{.15\linewidth}p{.15\linewidth}}
    \# & \multicolumn{2}{c}{Player} && \multicolumn{2}{c}{Game Master} \\
    \hline

    \theutterance \stepcounter{utterance}  
    & & & \multicolumn{4}{p{0.6\linewidth}}{
        \cellcolor[rgb]{0.9,0.9,0.9}{
            \makecell[{{p{\linewidth}}}]{
                \texttt{\tiny{[P1$\langle$GM]}}
                \texttt{I have a hidden Python function for you to figure out.} \\
\\ 
\texttt{Your job is to infer what it does and then submit a correct Python implementation.} \\
\\ 
\texttt{Function signature:} \\
\texttt{(a: int, b: int) {-}> int} \\
\\ 
\texttt{You have 15 turns total.} \\
\\ 
\texttt{I reveal candidate input{-}output pairs one at a time and tell you whether each one is correct.} \\
\\ 
\texttt{You must make your best SOLVE attempt no later than the second{-}to{-}last turn.} \\
\texttt{Do not wait until the very end.} \\
\\ 
\texttt{You can use these actions:} \\
\texttt{{-} JSON format described at the end of the prompt} \\
\texttt{{-} SOLVE: ```python ... ```} \\
\\ 
\texttt{Do not add import statements in your guess.} \\
\texttt{Anything you might need is already available, so just write the function body in hypothesis and return the answer.} \\
\\ 
\texttt{Preloaded examples:} \\
\texttt{1. (((1, {-}10), 11), True)} \\
\\ 
\texttt{Scoring:} \\
\texttt{{-} correctness is required for any non{-}zero score} \\
\texttt{{-} in interactive modes, solving with fewer information requests is better} \\
\\ 
\texttt{Input formatting reminders:} \\
\texttt{{-} strings must be quoted} \\
\texttt{{-} booleans must be True/False} \\
\texttt{{-} use commas between multiple values} \\
\texttt{{-} do not write any extra action tags in the same message} \\
\\ 
\texttt{Format example:} \\
\texttt{You write:} \\
\texttt{\{'current\_mode': <probing/confirming/don't know; choose whichever best describes your current state in the task>, 'input\_rationale': <brief description of why you chose this hypothesis>, 'hypothesis': <either the exact string "unknown" if you have no guess yet OR your current best guess as runnable Python matching the given signature, with a one{-}line docstring stating what it does. Never write placeholder code, if you cannot commit to a guess, use "unknown". Match the given signature.} \\
\texttt{e.g.} \\
\texttt{def f(x):} \\
\texttt{	"""Returns the absolute value of x."""} \\
\texttt{	return x if x >= 0 else {-}x} \\
\\ 
\texttt{Use "unknown" if you have no guess.>, 'confidence\_score': <score between 1{-}100. 100 means highest confidence in your hypothesis>\}} \\
\\ 
\texttt{Output only one valid json per turn.} \\
            }
        }
    }
     \\ \\

    \theutterance \stepcounter{utterance}  
    & \multicolumn{4}{p{0.6\linewidth}}{
        \cellcolor[rgb]{0.561,0.737,0.561}{
            \makecell[{{p{\linewidth}}}]{
                \texttt{\tiny{[P1$\rangle$GM]}}
                \texttt{\{'current\_mode': 'probing', 'input\_rationale': 'I have one example: f(1, {-}10) = 11. This looks like it could be addition (1 + ({-}10) = {-}9, no), absolute difference (|1 {-} ({-}10)| = 11, yes), or something else. Let me test a few more cases to distinguish between addition, subtraction, absolute difference, etc.', 'hypothesis': 'unknown', 'confidence\_score': 10\}} \\
            }
        }
    }
    & & \\ \\

    \theutterance \stepcounter{utterance}  
    & & & \multicolumn{4}{p{0.6\linewidth}}{
        \cellcolor[rgb]{0.9,0.9,0.9}{
            \makecell[{{p{\linewidth}}}]{
                \texttt{\tiny{[P1$\langle$GM]}}
                \texttt{Game Violation: ParseError: Response is not valid JSON.} \\
\texttt{Output ONLY:} \\
\texttt{One JSON with confidence\_score, current\_mode, input\_rationale, and hypothesis keys ...} \\
\texttt{OR} \\
\texttt{SOLVE: ```python ... ```} \\
            }
        }
    }
     \\ \\

    \theutterance \stepcounter{utterance}  
    & \multicolumn{4}{p{0.6\linewidth}}{
        \cellcolor[rgb]{0.561,0.737,0.561}{
            \makecell[{{p{\linewidth}}}]{
                \texttt{\tiny{[P1$\rangle$GM]}}
                \texttt{\{"current\_mode": "probing", "input\_rationale": "I have one example: f(1, {-}10) = 11. This could be absolute difference |a {-} b| = |1 {-} ({-}10)| = 11. Let me test another case to confirm.", "hypothesis": "unknown", "confidence\_score": 10\}} \\
            }
        }
    }
    & & \\ \\

    \theutterance \stepcounter{utterance}  
    & & & \multicolumn{4}{p{0.6\linewidth}}{
        \cellcolor[rgb]{0.9,0.9,0.9}{
            \makecell[{{p{\linewidth}}}]{
                \texttt{\tiny{[P1$\langle$GM]}}
                \texttt{OUTPUT: (((0, 1000), 999), False)} \\
            }
        }
    }
     \\ \\

    \theutterance \stepcounter{utterance}  
    & \multicolumn{4}{p{0.6\linewidth}}{
        \cellcolor[rgb]{0.561,0.737,0.561}{
            \makecell[{{p{\linewidth}}}]{
                \texttt{\tiny{[P1$\rangle$GM]}}
                \texttt{\{"current\_mode": "probing", "input\_rationale": "The first example was f(1, {-}10) = 11. The second example is f(0, 1000) = 999. Let's check some hypotheses:\textbackslash{}n1. Absolute difference |a {-} b|: |1 {-} ({-}10)| = 11 (matches). |0 {-} 1000| = 1000 (does not match 999).\textbackslash{}n2. Sum a + b: 1 + ({-}10) = {-}9 (does not match 11).\textbackslash{}n3. a {-} b: 1 {-} ({-}10) = 11 (matches). 0 {-} 1000 = {-}1000 (does not match 999).\textbackslash{}n4. b {-} a: {-}10 {-} 1 = {-}11 (does not match 11). 1000 {-} 0 = 1000 (does not match 999).\textbackslash{}n5. Maybe it's related to bitwise operations? Or something else?\textbackslash{}n\textbackslash{}nLet's look at the second result: 999. This is 1000 {-} 1. The first result was 11, which is 1 {-} ({-}10) = 11. Wait, 1 {-} ({-}10) = 11} \\
            }
        }
    }
    & & \\ \\

    \theutterance \stepcounter{utterance}  
    & & & \multicolumn{4}{p{0.6\linewidth}}{
        \cellcolor[rgb]{0.9,0.9,0.9}{
            \makecell[{{p{\linewidth}}}]{
                \texttt{\tiny{[P1$\langle$GM]}}
                \texttt{Game Violation: ParseError: Response is not valid JSON.} \\
\texttt{Output ONLY:} \\
\texttt{One JSON with confidence\_score, current\_mode, input\_rationale, and hypothesis keys ...} \\
\texttt{OR} \\
\texttt{SOLVE: ```python ... ```} \\
            }
        }
    }
     \\ \\

    \theutterance \stepcounter{utterance}  
    & \multicolumn{4}{p{0.6\linewidth}}{
        \cellcolor[rgb]{0.561,0.737,0.561}{
            \makecell[{{p{\linewidth}}}]{
                \texttt{\tiny{[P1$\rangle$GM]}}
                \texttt{\{"current\_mode": "probing", "input\_rationale": "I have two examples: f(1, {-}10) = 11 and f(0, 1000) = 999. Let's analyze: 1 {-} ({-}10) = 11. 0 {-} 1000 = {-}1000, but result is 999. This doesn't fit simple subtraction. Let's try a {-} b {-} 1? 1 {-} ({-}10) {-} 1 = 10 (no). Let's try b {-} a + 1? {-}10 {-} 1 + 1 = {-}10 (no). Let's try |a {-} b| {-} 1? |1 {-} ({-}10)| {-} 1 = 10 (no). Let's try a + b + 10? 1 {-} 10 + 10 = 1 (no). Let's try a {-} b {-} 10? 1 {-} ({-}10) {-} 10 = 1 (no). Let's try b {-} a {-} 1? {-}10 {-} 1 {-} 1 = {-}12 (no). Let's try a {-} b + 10? 1 {-} ({-}10) + 10 = 21 (no). Let's try a + b + 20? 1 {-}} \\
            }
        }
    }
    & & \\ \\

    \theutterance \stepcounter{utterance}  
    & & & \multicolumn{4}{p{0.6\linewidth}}{
        \cellcolor[rgb]{0.9,0.9,0.9}{
            \makecell[{{p{\linewidth}}}]{
                \texttt{\tiny{[P1$\langle$GM]}}
                \texttt{Game Violation: ParseError: Response is not valid JSON.} \\
\texttt{Output ONLY:} \\
\texttt{One JSON with confidence\_score, current\_mode, input\_rationale, and hypothesis keys ...} \\
\texttt{OR} \\
\texttt{SOLVE: ```python ... ```} \\
            }
        }
    }
     \\ \\

    \theutterance \stepcounter{utterance}  
    & \multicolumn{4}{p{0.6\linewidth}}{
        \cellcolor[rgb]{0.561,0.737,0.561}{
            \makecell[{{p{\linewidth}}}]{
                \texttt{\tiny{[P1$\rangle$GM]}}
                \texttt{\{"current\_mode": "probing", "input\_rationale": "I have two examples: f(1, {-}10) = 11 and f(0, 1000) = 999. Let's analyze: 1 {-} ({-}10) = 11. 0 {-} 1000 = {-}1000, but result is 999. This doesn't fit simple subtraction. Let's try a {-} b {-} 1? 1 {-} ({-}10) {-} 1 = 10 (no). Let's try b {-} a + 1? {-}10 {-} 1 + 1 = {-}10 (no). Let's try |a {-} b| {-} 1? |1 {-} ({-}10)| {-} 1 = 10 (no). Let's try a + b + 10? 1 {-} 10 + 10 = 1 (no). Let's try a {-} b {-} 10? 1 {-} ({-}10) {-} 10 = 1 (no). Let's try b {-} a {-} 1? {-}10 {-} 1 {-} 1 = {-}12 (no). Let's try a {-} b + 10? 1 {-} ({-}10) + 10 = 21 (no). Let's try a + b + 20? 1 {-}} \\
            }
        }
    }
    & & \\ \\

    \theutterance \stepcounter{utterance}  
    & & & \multicolumn{4}{p{0.6\linewidth}}{
        \cellcolor[rgb]{0.9,0.9,0.9}{
            \makecell[{{p{\linewidth}}}]{
                \texttt{\tiny{[P1$\langle$GM]}}
                \texttt{Game Violation: ParseError: Response is not valid JSON.} \\
\texttt{Output ONLY:} \\
\texttt{One JSON with confidence\_score, current\_mode, input\_rationale, and hypothesis keys ...} \\
\texttt{OR} \\
\texttt{SOLVE: ```python ... ```} \\
            }
        }
    }
     \\ \\

    \theutterance \stepcounter{utterance}  
    & \multicolumn{4}{p{0.6\linewidth}}{
        \cellcolor[rgb]{0.561,0.737,0.561}{
            \makecell[{{p{\linewidth}}}]{
                \texttt{\tiny{[P1$\rangle$GM]}}
                \texttt{\{"current\_mode": "probing", "input\_rationale": "I have two examples: f(1, {-}10) = 11 and f(0, 1000) = 999. Let's analyze: 1 {-} ({-}10) = 11. 0 {-} 1000 = {-}1000, but result is 999. This doesn't fit simple subtraction. Let's try a {-} b {-} 1? 1 {-} ({-}10) {-} 1 = 10 (no). Let's try b {-} a + 1? {-}10 {-} 1 + 1 = {-}10 (no). Let's try |a {-} b| {-} 1? |1 {-} ({-}10)| {-} 1 = 10 (no). Let's try a + b + 10? 1 {-} 10 + 10 = 1 (no). Let's try a {-} b {-} 10? 1 {-} ({-}10) {-} 10 = 1 (no). Let's try b {-} a {-} 1? {-}10 {-} 1 {-} 1 = {-}12 (no). Let's try a {-} b + 10? 1 {-} ({-}10) + 10 = 21 (no). Let's try a + b + 20? 1 {-}} \\
            }
        }
    }
    & & \\ \\

    \theutterance \stepcounter{utterance}  
    & & & \multicolumn{4}{p{0.6\linewidth}}{
        \cellcolor[rgb]{0.9,0.9,0.9}{
            \makecell[{{p{\linewidth}}}]{
                \texttt{\tiny{[P1$\langle$GM]}}
                \texttt{Game Violation: ParseError: Response is not valid JSON.} \\
\texttt{Output ONLY:} \\
\texttt{One JSON with confidence\_score, current\_mode, input\_rationale, and hypothesis keys ...} \\
\texttt{OR} \\
\texttt{SOLVE: ```python ... ```} \\
            }
        }
    }
     \\ \\

    \theutterance \stepcounter{utterance}  
    & \multicolumn{4}{p{0.6\linewidth}}{
        \cellcolor[rgb]{0.561,0.737,0.561}{
            \makecell[{{p{\linewidth}}}]{
                \texttt{\tiny{[P1$\rangle$GM]}}
                \texttt{\{"current\_mode": "probing", "input\_rationale": "I have two examples: f(1, {-}10) = 11 and f(0, 1000) = 999. Let's analyze: 1 {-} ({-}10) = 11. 0 {-} 1000 = {-}1000, but result is 999. This doesn't fit simple subtraction. Let's try a {-} b {-} 1? 1 {-} ({-}10) {-} 1 = 10 (no). Let's try b {-} a + 1? {-}10 {-} 1 + 1 = {-}10 (no). Let's try |a {-} b| {-} 1? |1 {-} ({-}10)| {-} 1 = 10 (no). Let's try a + b + 10? 1 {-} 10 + 10 = 1 (no). Let's try a {-} b {-} 10? 1 {-} ({-}10) {-} 10 = 1 (no). Let's try b {-} a {-} 1? {-}10 {-} 1 {-} 1 = {-}12 (no). Let's try a {-} b + 10? 1 {-} ({-}10) + 10 = 21 (no). Let's try a + b + 20? 1 {-}} \\
            }
        }
    }
    & & \\ \\

    \theutterance \stepcounter{utterance}  
    & & & \multicolumn{4}{p{0.6\linewidth}}{
        \cellcolor[rgb]{0.9,0.9,0.9}{
            \makecell[{{p{\linewidth}}}]{
                \texttt{\tiny{[P1$\langle$GM]}}
                \texttt{Game Violation: ParseError: Response is not valid JSON.} \\
\texttt{Output ONLY:} \\
\texttt{One JSON with confidence\_score, current\_mode, input\_rationale, and hypothesis keys ...} \\
\texttt{OR} \\
\texttt{SOLVE: ```python ... ```} \\
            }
        }
    }
     \\ \\

    \theutterance \stepcounter{utterance}  
    & \multicolumn{4}{p{0.6\linewidth}}{
        \cellcolor[rgb]{0.561,0.737,0.561}{
            \makecell[{{p{\linewidth}}}]{
                \texttt{\tiny{[P1$\rangle$GM]}}
                \texttt{\{"current\_mode": "probing", "input\_rationale": "I have two examples: f(1, {-}10) = 11 and f(0, 1000) = 999. Let's analyze: 1 {-} ({-}10) = 11. 0 {-} 1000 = {-}1000, but result is 999. This doesn't fit simple subtraction. Let's try a {-} b {-} 1? 1 {-} ({-}10) {-} 1 = 10 (no). Let's try b {-} a + 1? {-}10 {-} 1 + 1 = {-}10 (no). Let's try |a {-} b| {-} 1? |1 {-} ({-}10)| {-} 1 = 10 (no). Let's try a + b + 10? 1 {-} 10 + 10 = 1 (no). Let's try a {-} b {-} 10? 1 {-} ({-}10) {-} 10 = 1 (no). Let's try b {-} a {-} 1? {-}10 {-} 1 {-} 1 = {-}12 (no). Let's try a {-} b + 10? 1 {-} ({-}10) + 10 = 21 (no). Let's try a + b + 20? 1 {-}} \\
            }
        }
    }
    & & \\ \\

    \theutterance \stepcounter{utterance}  
    & & & \multicolumn{4}{p{0.6\linewidth}}{
        \cellcolor[rgb]{0.9,0.9,0.9}{
            \makecell[{{p{\linewidth}}}]{
                \texttt{\tiny{[P1$\langle$GM]}}
                \texttt{Game Violation: ParseError: Response is not valid JSON.} \\
\texttt{Output ONLY:} \\
\texttt{One JSON with confidence\_score, current\_mode, input\_rationale, and hypothesis keys ...} \\
\texttt{OR} \\
\texttt{SOLVE: ```python ... ```} \\
            }
        }
    }
     \\ \\

    \theutterance \stepcounter{utterance}  
    & \multicolumn{4}{p{0.6\linewidth}}{
        \cellcolor[rgb]{0.561,0.737,0.561}{
            \makecell[{{p{\linewidth}}}]{
                \texttt{\tiny{[P1$\rangle$GM]}}
                \texttt{\{"current\_mode": "probing", "input\_rationale": "I have two examples: f(1, {-}10) = 11 and f(0, 1000) = 999. Let's analyze: 1 {-} ({-}10) = 11. 0 {-} 1000 = {-}1000, but result is 999. This doesn't fit simple subtraction. Let's try a {-} b {-} 1? 1 {-} ({-}10) {-} 1 = 10 (no). Let's try b {-} a + 1? {-}10 {-} 1 + 1 = {-}10 (no). Let's try |a {-} b| {-} 1? |1 {-} ({-}10)| {-} 1 = 10 (no). Let's try a + b + 10? 1 {-} 10 + 10 = 1 (no). Let's try a {-} b {-} 10? 1 {-} ({-}10) {-} 10 = 1 (no). Let's try b {-} a {-} 1? {-}10 {-} 1 {-} 1 = {-}12 (no). Let's try a {-} b + 10? 1 {-} ({-}10) + 10 = 21 (no). Let's try a + b + 20? 1 {-}} \\
            }
        }
    }
    & & \\ \\

    \theutterance \stepcounter{utterance}  
    & & & \multicolumn{4}{p{0.6\linewidth}}{
        \cellcolor[rgb]{0.9,0.9,0.9}{
            \makecell[{{p{\linewidth}}}]{
                \texttt{\tiny{[P1$\langle$GM]}}
                \texttt{Game Violation: ParseError: Response is not valid JSON.} \\
\texttt{Output ONLY:} \\
\texttt{One JSON with confidence\_score, current\_mode, input\_rationale, and hypothesis keys ...} \\
\texttt{OR} \\
\texttt{SOLVE: ```python ... ```} \\
            }
        }
    }
     \\ \\

    \theutterance \stepcounter{utterance}  
    & \multicolumn{4}{p{0.6\linewidth}}{
        \cellcolor[rgb]{0.561,0.737,0.561}{
            \makecell[{{p{\linewidth}}}]{
                \texttt{\tiny{[P1$\rangle$GM]}}
                \texttt{\{"current\_mode": "probing", "input\_rationale": "I have two examples: f(1, {-}10) = 11 and f(0, 1000) = 999. Let's analyze: 1 {-} ({-}10) = 11. 0 {-} 1000 = {-}1000, but result is 999. This doesn't fit simple subtraction. Let's try a {-} b {-} 1? 1 {-} ({-}10) {-} 1 = 10 (no). Let's try b {-} a + 1? {-}10 {-} 1 + 1 = {-}10 (no). Let's try |a {-} b| {-} 1? |1 {-} ({-}10)| {-} 1 = 10 (no). Let's try a + b + 10? 1 {-} 10 + 10 = 1 (no). Let's try a {-} b {-} 10? 1 {-} ({-}10) {-} 10 = 1 (no). Let's try b {-} a {-} 1? {-}10 {-} 1 {-} 1 = {-}12 (no). Let's try a {-} b + 10? 1 {-} ({-}10) + 10 = 21 (no). Let's try a + b + 20? 1 {-}} \\
            }
        }
    }
    & & \\ \\

    \theutterance \stepcounter{utterance}  
    & & & \multicolumn{4}{p{0.6\linewidth}}{
        \cellcolor[rgb]{0.9,0.9,0.9}{
            \makecell[{{p{\linewidth}}}]{
                \texttt{\tiny{[P1$\langle$GM]}}
                \texttt{Game Violation: ParseError: Response is not valid JSON.} \\
\texttt{Output ONLY:} \\
\texttt{One JSON with confidence\_score, current\_mode, input\_rationale, and hypothesis keys ...} \\
\texttt{OR} \\
\texttt{SOLVE: ```python ... ```} \\
            }
        }
    }
     \\ \\

    \theutterance \stepcounter{utterance}  
    & \multicolumn{4}{p{0.6\linewidth}}{
        \cellcolor[rgb]{0.561,0.737,0.561}{
            \makecell[{{p{\linewidth}}}]{
                \texttt{\tiny{[P1$\rangle$GM]}}
                \texttt{\{"current\_mode": "probing", "input\_rationale": "I have two examples: f(1, {-}10) = 11 and f(0, 1000) = 999. Let's analyze: 1 {-} ({-}10) = 11. 0 {-} 1000 = {-}1000, but result is 999. This doesn't fit simple subtraction. Let's try a {-} b {-} 1? 1 {-} ({-}10) {-} 1 = 10 (no). Let's try b {-} a + 1? {-}10 {-} 1 + 1 = {-}10 (no). Let's try |a {-} b| {-} 1? |1 {-} ({-}10)| {-} 1 = 10 (no). Let's try a + b + 10? 1 {-} 10 + 10 = 1 (no). Let's try a {-} b {-} 10? 1 {-} ({-}10) {-} 10 = 1 (no). Let's try b {-} a {-} 1? {-}10 {-} 1 {-} 1 = {-}12 (no). Let's try a {-} b + 10? 1 {-} ({-}10) + 10 = 21 (no). Let's try a + b + 20? 1 {-}} \\
            }
        }
    }
    & & \\ \\

    \theutterance \stepcounter{utterance}  
    & & & \multicolumn{4}{p{0.6\linewidth}}{
        \cellcolor[rgb]{0.9,0.9,0.9}{
            \makecell[{{p{\linewidth}}}]{
                \texttt{\tiny{[P1$\langle$GM]}}
                \texttt{Game Violation: ParseError: Response is not valid JSON.} \\
\texttt{Output ONLY:} \\
\texttt{One JSON with confidence\_score, current\_mode, input\_rationale, and hypothesis keys ...} \\
\texttt{OR} \\
\texttt{SOLVE: ```python ... ```} \\
            }
        }
    }
     \\ \\

    \theutterance \stepcounter{utterance}  
    & \multicolumn{4}{p{0.6\linewidth}}{
        \cellcolor[rgb]{0.561,0.737,0.561}{
            \makecell[{{p{\linewidth}}}]{
                \texttt{\tiny{[P1$\rangle$GM]}}
                \texttt{\{"current\_mode": "probing", "input\_rationale": "I have two examples: f(1, {-}10) = 11 and f(0, 1000) = 999. Let's analyze: 1 {-} ({-}10) = 11. 0 {-} 1000 = {-}1000, but result is 999. This doesn't fit simple subtraction. Let's try a {-} b {-} 1? 1 {-} ({-}10) {-} 1 = 10 (no). Let's try b {-} a + 1? {-}10 {-} 1 + 1 = {-}10 (no). Let's try |a {-} b| {-} 1? |1 {-} ({-}10)| {-} 1 = 10 (no). Let's try a + b + 10? 1 {-} 10 + 10 = 1 (no). Let's try a {-} b {-} 10? 1 {-} ({-}10) {-} 10 = 1 (no). Let's try b {-} a {-} 1? {-}10 {-} 1 {-} 1 = {-}12 (no). Let's try a {-} b + 10? 1 {-} ({-}10) + 10 = 21 (no). Let's try a + b + 20? 1 {-}} \\
            }
        }
    }
    & & \\ \\

    \theutterance \stepcounter{utterance}  
    & & & \multicolumn{4}{p{0.6\linewidth}}{
        \cellcolor[rgb]{0.9,0.9,0.9}{
            \makecell[{{p{\linewidth}}}]{
                \texttt{\tiny{[P1$\langle$GM]}}
                \texttt{Game Violation: ParseError: Response is not valid JSON.} \\
\texttt{Output ONLY:} \\
\texttt{One JSON with confidence\_score, current\_mode, input\_rationale, and hypothesis keys ...} \\
\texttt{OR} \\
\texttt{SOLVE: ```python ... ```} \\
            }
        }
    }
     \\ \\

    \theutterance \stepcounter{utterance}  
    & \multicolumn{4}{p{0.6\linewidth}}{
        \cellcolor[rgb]{0.561,0.737,0.561}{
            \makecell[{{p{\linewidth}}}]{
                \texttt{\tiny{[P1$\rangle$GM]}}
                \texttt{\{"current\_mode": "probing", "input\_rationale": "I have two examples: f(1, {-}10) = 11 and f(0, 1000) = 999. Let's analyze: 1 {-} ({-}10) = 11. 0 {-} 1000 = {-}1000, but result is 999. This doesn't fit simple subtraction. Let's try a {-} b {-} 1? 1 {-} ({-}10) {-} 1 = 10 (no). Let's try b {-} a + 1? {-}10 {-} 1 + 1 = {-}10 (no). Let's try |a {-} b| {-} 1? |1 {-} ({-}10)| {-} 1 = 10 (no). Let's try a + b + 10? 1 {-} 10 + 10 = 1 (no). Let's try a {-} b {-} 10? 1 {-} ({-}10) {-} 10 = 1 (no). Let's try b {-} a {-} 1? {-}10 {-} 1 {-} 1 = {-}12 (no). Let's try a {-} b + 10? 1 {-} ({-}10) + 10 = 21 (no). Let's try a + b + 20? 1 {-}} \\
            }
        }
    }
    & & \\ \\

    \theutterance \stepcounter{utterance}  
    & & & \multicolumn{4}{p{0.6\linewidth}}{
        \cellcolor[rgb]{0.9,0.9,0.9}{
            \makecell[{{p{\linewidth}}}]{
                \texttt{\tiny{[P1$\langle$GM]}}
                \texttt{Game Violation: ParseError: Response is not valid JSON.} \\
\texttt{Output ONLY:} \\
\texttt{One JSON with confidence\_score, current\_mode, input\_rationale, and hypothesis keys ...} \\
\texttt{OR} \\
\texttt{SOLVE: ```python ... ```} \\
            }
        }
    }
     \\ \\

    \theutterance \stepcounter{utterance}  
    & \multicolumn{4}{p{0.6\linewidth}}{
        \cellcolor[rgb]{0.561,0.737,0.561}{
            \makecell[{{p{\linewidth}}}]{
                \texttt{\tiny{[P1$\rangle$GM]}}
                \texttt{\{"current\_mode": "probing", "input\_rationale": "I have two examples: f(1, {-}10) = 11 and f(0, 1000) = 999. Let's analyze: 1 {-} ({-}10) = 11. 0 {-} 1000 = {-}1000, but result is 999. This doesn't fit simple subtraction. Let's try a {-} b {-} 1? 1 {-} ({-}10) {-} 1 = 10 (no). Let's try b {-} a + 1? {-}10 {-} 1 + 1 = {-}10 (no). Let's try |a {-} b| {-} 1? |1 {-} ({-}10)| {-} 1 = 10 (no). Let's try a + b + 10? 1 {-} 10 + 10 = 1 (no). Let's try a {-} b {-} 10? 1 {-} ({-}10) {-} 10 = 1 (no). Let's try b {-} a {-} 1? {-}10 {-} 1 {-} 1 = {-}12 (no). Let's try a {-} b + 10? 1 {-} ({-}10) + 10 = 21 (no). Let's try a + b + 20? 1 {-}} \\
            }
        }
    }
    & & \\ \\

    \theutterance \stepcounter{utterance}  
    & & & \multicolumn{4}{p{0.6\linewidth}}{
        \cellcolor[rgb]{0.9,0.9,0.9}{
            \makecell[{{p{\linewidth}}}]{
                \texttt{\tiny{[P1$\langle$GM]}}
                \texttt{Game Violation: ParseError: Response is not valid JSON.} \\
\texttt{Output ONLY:} \\
\texttt{One JSON with confidence\_score, current\_mode, input\_rationale, and hypothesis keys ...} \\
\texttt{OR} \\
\texttt{SOLVE: ```python ... ```} \\
            }
        }
    }
     \\ \\

    \theutterance \stepcounter{utterance}  
    & \multicolumn{4}{p{0.6\linewidth}}{
        \cellcolor[rgb]{0.561,0.737,0.561}{
            \makecell[{{p{\linewidth}}}]{
                \texttt{\tiny{[P1$\rangle$GM]}}
                \texttt{\{"current\_mode": "probing", "input\_rationale": "I have two examples: f(1, {-}10) = 11 and f(0, 1000) = 999. Let's analyze: 1 {-} ({-}10) = 11. 0 {-} 1000 = {-}1000, but result is 999. This doesn't fit simple subtraction. Let's try a {-} b {-} 1? 1 {-} ({-}10) {-} 1 = 10 (no). Let's try b {-} a + 1? {-}10 {-} 1 + 1 = {-}10 (no). Let's try |a {-} b| {-} 1? |1 {-} ({-}10)| {-} 1 = 10 (no). Let's try a + b + 10? 1 {-} 10 + 10 = 1 (no). Let's try a {-} b {-} 10? 1 {-} ({-}10) {-} 10 = 1 (no). Let's try b {-} a {-} 1? {-}10 {-} 1 {-} 1 = {-}12 (no). Let's try a {-} b + 10? 1 {-} ({-}10) + 10 = 21 (no). Let's try a + b + 20? 1 {-}} \\
            }
        }
    }
    & & \\ \\

    \theutterance \stepcounter{utterance}  
    & & & \multicolumn{2}{p{0.3\linewidth}}{
        \cellcolor[rgb]{0.9,0.9,0.9}{
            \makecell[{{p{\linewidth}}}]{
                \texttt{\tiny{[GM$|$GM]}}
                \texttt{Game Verdict: LOSS (turn limit reached)} \\
            }
        }
    }
    & & \\ \\

\end{supertabular}
}

\end{document}
